\section{Introduction}
\label{sec:intro}

Unsupervised environment design (UED) trains a generalist by choosing which levels it
practises on. Most current methods are motivated by minimax regret: a level is worth
replaying when the gap between the best achievable return on it and the agent's return is
large~\citep{jiang2021robustplr,dennis2020paired}. That gap needs the optimal return on
every level, which is unknown in almost every domain of interest. So the methods use
proxies computed from the agent's own rollouts. Robust Prioritised Level Replay
(\PLRp{})~\citep{jiang2021robustplr}, which builds on level replay~\citep{jiang2021plr},
ranks levels by positive value loss (PVL) or by MaxMC, the gap between the best return
seen on a level and the critic's value. Newer regret proxies include maximised negative
advantage (MNA)~\citep{mead2026mna} and a transition-prediction term added to PVL in
TRACED~\citep{cho2026traced}. Other scores target something else. Sampling for
Learnability (SFL) ranks levels by $p(1-p)$, where $p$ is the success
rate~\citep{rutherford2024noregrets}; No Regrets proposes it in place of regret
approximations, as the probability of improving times the room left to improve, not as
an estimate of regret.
% src: docs/paper/sources-sfl.md section 1.1 ("justified as (probability of improving) x
%      (room left to improve)")
PACE~\citep{yuan2026pace} scores the parameter change of a provisional update, a measure
of learning progress, and the spread of a critic ensemble, an option in the minimax
library~\citep{minimax2023code}, has antecedents in curricula driven by value-ensemble
disagreement~\citep{zhang2020vds,kanitscheider2022autodime}.
% src: docs/paper/sources-recent.md sections (b2) and (d)

\paragraph{The problem} The regret proxies are treated as regret estimates, but whether
they rank levels by regret has not been checked where it matters. For the scores that
target something else, how far they agree with regret is a comparison of targets, which
is also unmeasured. The closest test, in No
Regrets~\citep{rutherford2024noregrets}, correlates PVL and MaxMC with \emph{success
rate} for one policy, one seed and one checkpoint in single-agent JaxNav, on levels
stratified by that same policy's success; it reports Pearson $r = 0.61$ for MaxMC and
$0.51$ for PVL against success,
% src: docs/paper/sources-sfl.md section 3.1 (Fig. 2a, 2b) and section 3.3
and never compares any score with regret. Exact-regret curricula exist in gridworlds,
where a shortest-path solver gives the optimum~\citep{beukman2024remidi,abdelsadek2025goal},
but they use the oracle to sample levels rather than to audit the estimators. The newer
scores are judged by the test performance of the curriculum they
drive~\citep{mead2026mna,cho2026traced,yuan2026pace}, which mixes the quality of the score
with everything else the curriculum does. In our literature search we found no paper from
2024--26 that compares UED level scores with exact or best-known regret across training.
% src: docs/plan/two-track-plan.md section 2.7 ("Verdict")
\VERIFY{re-run the search before submission: ICLR 2027 submissions and NeurIPS 2027
accepted papers, as scheduled in docs/plan/two-track-plan.md section 2.7}

The difficulty is the answer key. Regret is defined in the agent's own units, its
discounted return, so the reference has to be the optimum of that same objective, on the
same dynamics, at the same control rate and under the same termination rules. A lap-time
optimum or a success indicator measures something else. Any mismatch between the
reference problem and the agent's problem then shows up as ``regret'' that no policy
could remove.

\paragraph{Our instrument} We build a continuous-control domain in which the level-wise
optimum of the agent's own objective can be computed (\cref{sec:env,sec:answer}). A level
is an open, procedurally generated race-track strip, a car, a hidden grip value and a
start state. One vehicle model, a dynamic single-track car in curvilinear coordinates,
is written once and compiled twice: to JAX for the reinforcement-learning simulator and to
CasADi for optimal control. On $10^{4}$ random states the two agree to
% src: tests/test_dynamics.py::test_jax_casadi_parity (threshold 1e-5);
%      paper/tools/gate_values.py (measured 1.65e-14)
$2\times10^{-14}$ (\cref{tab:gates}). On each level a
\emph{training-objective OCP} (TO-OCP) maximises the agent's discounted reward over the
agent's 60\,s window. It uses the environment's RK2 integrator with ten 5\,ms substeps per
50\,ms control step, a zero-order hold at 20\,Hz, the environment's termination rules as
constraints (tightened slightly where they differ, which can only lower the optimum), and
the agent's reward as its objective. We replay the optimal controls
open-loop in the float64 simulator and accept a label only if the replayed return matches.
The solver finds local optima, so we report its value $\Jstar$ as a \emph{best-known}
optimum and the measured regret $\Jstar - \Jpi$ as a \emph{lower bound} on the true
regret (\cref{sec:answer-lb}).

\paragraph{What we measure} We ask three questions (\cref{sec:protocol}). \textbf{RQ1,
fidelity:} how well do PVL, L1 value loss, MaxMC (as implemented and in a unit-consistent
variant), MNA and ensemble disagreement rank levels by regret, on fresh rollouts and
inside the buffers that real \PLRp{} and SFL runs keep? We measure SFL's learnability
the same way, but because it targets learning potential rather than regret, its
agreement with regret is a comparison of two targets, not a test of an
estimator. \textbf{RQ2,
dynamics:} how does fidelity change with training phase, critic privilege and
observability, and which measured term explains the change? Two exact identities split
each score's error into critic bias, dispersion, the gap between the best observed and
the optimal return, success saturation and staleness (\cref{sec:est-identities}).
\textbf{RQ3, consequence:} on one labelled pool, does a \PLRp{} curriculum scored by
oracle regret beat \PLRp{} with PVL or MaxMC, SFL and domain randomisation on held-out
regret? A replication in JaxUED's maze~\citep{coward2024jaxued}, run with JaxUED's own
code and an exact breadth-first-search optimum, checks that the racing results are not an
artefact of one bespoke domain. The scorer, the window design, the level partition, the
endpoints and the decision rules are frozen in two tags, \code{audit-v1} and
\code{protocol-v1}, before any full-scale result exists (\cref{sec:prereg}).

\paragraph{Contributions}
% src: docs/plan/two-track-plan.md section 2.1 (contributions 1-4)
\begin{enumerate}
  \item A JAX racing benchmark of about \num{5400} levels, each labelled with the
        replay-verified best-known optimum of the agent's own objective.
  \item To our knowledge, the first audit of UED level scores against ground-truth (here
        best-known) regret across training, both on fresh rollouts and inside real
        curriculum buffers.
  \item A comparison of curricula on a single shared labelled pool, including an
        oracle-regret arm.
  \item A replication of the audit in JaxUED's maze, using JaxUED's own training code at
        its defaults.
\end{enumerate}
The algebraic identities we use to decompose score error (\cref{sec:est-identities}) are
a measurement framework for RQ2, not a contribution.
% src: docs/plan/two-track-plan.md section 2.1 ("The algebraic identities in E2 ...")

\paragraph{Scope} Every claim is about the racing domain and, where stated, the maze. The
racing reward is dense; whether our findings carry over to sparse-reward domains is a
conjecture we do not test (\cref{sec:threats}). No audit or curriculum experiment uses a
real circuit. Real circuits enter the paper in three ways only: twelve development
circuits supply the generator's corner statistics; the development checks of
\cref{sec:env,sec:answer} (one specialist agent, a minimum-time replay gate and the
TO-OCP spike) ran on development circuits only; and \cref{app:snapping} audits public
human telemetry, which no claim uses. No number in the paper is computed on the twelve
lockbox circuits held out for later work (\cref{sec:threats}).
% src: docs/plan/two-track-plan.md sections 2.4 ("Real circuits are left out of the paper"),
%      2.6 (reward density stated as conjecture) and 4 (lockbox rules)
